\begin{exercise}{54}
Let $I$ be a finite category.
\begin{enumerate}[label=(\arabic*)]
\item Show that $I$ is filtered if it has a terminal object.
\item Show that there are examples where $I$ is filtered but does not have a terminal object.
\item Show that, if $I$ is a poset, it is filtered if and only if it has a terminal object.\footnote{The original says ``$I$ is a poset and filtered if and only if it has a terminal object''. The poset condition must instead be a hypothesis in (3). Indeed, the finite category generated by $f,g\colon0\rightrightarrows1$ and $h\colon1\to2$ with $hf=hg$ has terminal object $2$ but is not a poset, since $f\ne g$.}
\end{enumerate}
In particular, notice that this shows that there are many filtered categories which are not
posets.

\begin{proof}
(1) Extend any diagram $F\colon K\to I$ by sending the
new terminal object of $K^{\rhd}$ to a terminal object $t$ of $I$.
The unique maps $F(k)\to t$ are compatible with every morphism
of $K$, so this defines the required extension.

(2) Let $I$ be the one-object category whose endomorphism
monoid is $\{1,0\}$, with unit $1$ and absorbing element $0$.
For any $F\colon K\to I$, extend $F$ over $K^{\rhd}$ by
sending the new object to the sole object and every arrow
from $K$ to that object to $0$.
The compatibility equations are $0F(a)=0$, so this is a functor.
Thus $I$ is finite and filtered, but its sole object is not
terminal, since it has two endomorphisms.
It is also not a poset.

(3) If a finite poset $I$ is filtered, Exercise~\ref{ex:53}
provides an upper bound for the finite collection of all its
objects. This is a greatest element, hence a terminal object.
The converse is (1).
The example in (2) proves the final assertion; more generally,
any finite monoid with an absorbing element and more than one
element gives the same construction.
\end{proof}
\end{exercise}
